\newpage\EMPHASIZE{31. Typedefs}

\textsc{Objectives}
\begin{itemize}
\item Create typedef
\item Understand that typedef does not create a new type
\end{itemize}

I will show you how to create an alias for a type. This does not really
create a new type but simply gives a type another name.


\newpage\EMPHASIZE{typedef}

Just like we prefer to create constants instead of hardcoding constants:

\verb!const int MAX_AGE = 100;!

so we can also give type another (hopefully more readable) name. For
instance instead of

\begin{consolethree}[escapeinside=||]
const int CEO = 0;
const int MANAGER = 1;
const int FULLTIME = 2;
const int PARTTIME = 3;

int employeeCode = FULLTIME;
std::cin >> employeeCode;
\end{consolethree}

we can create a name for employee code type:

\begin{consolethree}[escapeinside=||]
typedef int EmployeeCode;

const EmployeeCode CEO = 0;
const EmployeeCode MANAGER = 1;
const EmployeeCode FULLTIME = 2;
const EmployeeCode PARTTIME = 3;

EmployeeCode employeeCode = FULLTIME;
std::cin >> employeeCode;
\end{consolethree}

The statement

\verb!typedef int EmployeeCode;!

basically tells C/C++ to replace \textit{EmployeeCode} with \textit{int}
before compiling the program. The format for a typedef looks like this:

\verb!typedef [type] [typedef name];!

\begin{ex}
Create a typedef for \textit{double} with the name
\textit{GPA}. Declare a variable \textit{johnDoeGPA} of type \textit{GPA} and
initialize it with the value of \textit{3.25}.
\end{ex}

It's important to remember that \textit{typedef}s are type
aliases. Why is this important? This means that you \textbf{cannot} have
the following overloaded functions like the following:

\begin{consolethree}[escapeinside=||]
#include <iostream>

typedef int EmployeeCode;

const EmployeeCode CEO = 0;
const EmployeeCode MANAGER = 1;
const EmployeeCode FULLTIME = 2;
const EmployeeCode PARTTIME = 3;

void print(int i)
{
    std::cout << i;
}

void print(EmployeeCode code)
{
    std::cout << "Employee code: " << code;
}

int main()
{
    print(42);
    print(CEO);
    return 0;
}
\end{consolethree}

Why? Because \textit{EmployeeCode} is just an alias for \textit{int}. In
other words as far as the compiler is concerned the code is the same as:

\begin{consolethree}[escapeinside=||]
...

void print(int i)
{
    std::cout << i;
}

void print(int code)
{
    std::cout << "Employee code: " << code;
}

...
\end{consolethree}

And of course you see the problem: There are two functions with the same
signatures and that's not allowed.

The following is an example on how to \textit{typedef} a reference type, a
constant type, and a pointer type (no surprises):

\begin{consolethree}[escapeinside=||]
typedef int & intr;   // typedef for int &
typedef const int cint; // typedef for const int
typedef int * pint;   // typedef for int *

int i = 42;
intr x = i;
cint j = 0;
pint p = &i;
\end{consolethree}

\begin{ex}
Make sure you create corresponding typedefs for
doubles. Test the typedefs.
\end{ex}


\newpage\EMPHASIZE{Typedef for Array Types}

You can also do typedef for array types. It's very
common to use an array of two integers to model a point in 2-dimensional
space.

\begin{consolethree}[escapeinside=||]
int p[2] = {3, 6}; // variable p models (3,6)
                   // in 2-dimensional space
\end{consolethree}

In the above code, we want to think of p as a point where \verb!p[0]! is
the x-coordinate of p while \verb!p[1]! is the y-coordinate of p.

You can do this:

\verb!typedef int Point[2];!

This is how you can use this \textit{typedef}:

\begin{consolethree}[escapeinside=||]
typedef int Point[2];
Point p = {3, 6};
\end{consolethree}

The format for creating typedefs for array types is this:

\verb!typedef [type] [array typedef][[size]];!

A very \textbf{common typo / error / ignorance / atrocity / etc}. is
this:

\verb!typedef int[2] Point;! // BADDDD!!!

Remember, the size of the array comes last.

\begin{ex}
Rewrite the following program using the above
\textit{typedef} for \textit{Point} and replace the type \verb!int[2]!
by Point.

\begin{consolethree}[escapeinside=||]
#include <iostream>

void print(int p[2])
{
    std::cout << "(" << p[0] << ", " << p[1] << ")";
}

void add(int sum[2], int p[2], int q[2])
{
    sum[0] = p[0] + q[0];
    sum[1] = p[1] + q[1];
}

int main()
{
    int q[2] = {3, 5};
    print(q);
    std::cout << std::endl;
    return 0;
}
\end{consolethree}
\end{ex}


\newpage\EMPHASIZE{Typedef For Pointer Types}

Try this:

\begin{consolethree}[escapeinside=||]
#include <iostream>

typedef int * IntPtr;

int main()
{
    IntPtr p = new int;
    *p = 42;
    std::cout << (*p) + 3 << std::endl;
    delete p;
}
\end{consolethree}

Get it?

\begin{ex}
Create a typedef for a pointer to doubles and use
this typedef in this code whenever possible. Fill in the missing code
and correct errors (yes there are errors).

\begin{consolethree}[escapeinside=||]
#include <iostream>

//---------------------------------------------------
// Allocate memory for p
//---------------------------------------------------
void constructArray(double * p, int size)
{
    p = new double[size];
}

//---------------------------------------------------
// Deallocate memory for p
//---------------------------------------------------
void destructArray(double * p)
{
    delete [] p;
}

//---------------------------------------------------
// Randomize the array p is pointing to.
//---------------------------------------------------
void randArray(double * p, int size)
{
    for (int i = 0; i < size; i++)
    {
        p[i] = rand() / RANDMAX;
    }
}

//---------------------------------------------------
// Print all the element p is pointing to
//---------------------------------------------------
void printArray(double * p, int size)
{
}

int main()
{
    srand();
    std::cout << "How many doubles do you want?";
    std::cin >> size;

    while (size > 0)
    {
        double * arr;
        constructArray(arr, size);
        randArray(p, size);
        printArray(p, size);
        destructArray(arr);
        std::cout << "How many doubles do you want?";
        std::cin >> size;
    }
}
\end{consolethree}
\end{ex}

\begin{ex}
Create a typedef for a \textbf{2}-dimensional
array in the following program and use it whenever possible.

\begin{consolethree}[escapeinside=||]
#include <iostream>

const int SIZE = 5;

// put your typedef here

void print(char x[SIZE][SIZE])
{
    for (int row = 0; row < SIZE; row++)
    {
        for (int col = 0; col < SIZE; col++)
        {
            std::cout << x[row][col];
        }
        std::cout << std::endl;
    }
}

int main()
{
    char a[SIZE][SIZE];

    int row = 0, col = 0;
    for (int i = 0; i < SIZE * SIZE; i++)
    {
        a[row][col] = char(i + 'a');
        col++;
        if (col == SIZE)
        {
            row++;
            col = 0;
        }
    }

    print(a);
    return 0;
}
\end{consolethree}
\end{ex}

\begin{ex}
Now rewrite the above code so that, without changing
the behavior of the program, the code has the following form:

\begin{consolethree}[escapeinside=||]
#include <iostream>

const int SIZE = 5;

// put your typedef here

void print(char x[SIZE][SIZE])
{
    for (int row = 0; row < SIZE; row++)
    {
        for (int col = 0; col < SIZE; col++)
        {
            std::cout << x[row][col];
        }
        std::cout << std::endl;
    }
}

int main()
{
    char a[SIZE][SIZE];

    for (int row = 0; row < SIZE; row++)
    {
        for (int col = 0; col < SIZE; col++)
        {
            a[row][col] = ____________;
        }
    }

    print(a);
    return 0;
}
\end{consolethree}
\end{ex}


\newpage\EMPHASIZE{typedef and Multi-file Compilation}

\textit{typedef}s can be placed in a header file.

When do you want to do that? When a \textit{typedef} is used by several
cpp files.

Not only that. If several functions are closely associated with the
\textit{typedef}, then the prototypes of these functions and the
\textit{typedef} should be in the same header file.

Here's an example:

\begin{consolethree}[escapeinside=||]
// employee.h
#ifndef EMPLOYEE_H
#define EMPLOYEE_H

#include <iostream>

typedef int EmployeeCode;

const EmployeeCode CEO = 0;
const EmployeeCode MANAGER = 1;
const EmployeeCode FULLTIME = 2;
const EmployeeCode PARTTIME = 3;

void print(EmployeeCode);
EmployeeCode promote(EmployeeCode);
EmployeeCode demote(EmployeeCode);

#endif EMPLOYEE_H
\end{consolethree}

\begin{consolethree}[escapeinside=||]
// employee.cpp
#include <iostream>
#include "employee.h"

void print(EmployeeCode code)
{
    std::cout << code;
}

EmployeeCode promote(EmployeeCode code)
{
    return (code > CEO ? code - 1 : code);
}

EmployeeCode demote(EmployeeCode)
{
    return (code < PARTTIME ? code + 1 : code);
}
\end{consolethree}

\begin{ex}
Rewrite this program so that you have two more file
for the array-related typedef and functions: a header file
\textit{Array.h} containing the typedef and function prototypes and C++
source file \textit{Array.cpp} for the definition of the array-related
functions. (See the previous section.)

\begin{consolethree}[escapeinside=||]
#include <iostream>

//---------------------------------------------------
// Allocate memory for p
//---------------------------------------------------
void constructArray(double * p, int size)
{
    p = new double[size];
}

//---------------------------------------------------
// Deallocate memory for p
//---------------------------------------------------
void destructArray(double * p)
{
    delete [] p;
}

//---------------------------------------------------
// Randomize the array p is pointing to with random
// doubles between 0.0 and 1.0.
//---------------------------------------------------
void randArray(double * p, int size)
{
    for (int i = 0; i < size; i++)
    {
        p[i] = rand() / RANDMAX;
    }
}

//---------------------------------------------------
// Print all the element p is pointing to
//---------------------------------------------------
void printArray(double * p, int size)
{
}

int main()
{
    srand(0);
    std::cout << "How many doubles do you want?";
    std::cin >> size;

    while (size > 0)
    {
        double * arr;
        constructArray(arr, size);
        randArray(p, size);
        printArray(p, size);
        destructArray(arr);
        std::cout << "How many doubles do you want?";
        std::cin >> size;
    }
}
\end{consolethree}
\end{ex}

\end{document}